% ===== §1 The three registers as coupled systems =====
\section{The Three Registers and the Emergence of Value}
\label{sec:registers}

\noindent
The phenomenological datum of the following section is the co-apprehension, by two persons, of a value generated between them. The adequate description of that datum requires a prior account of the structure within which value is located, and the present section supplies it. The account is Lacanian in its terms and dynamical in its construal: the three registers of the Real, the Symbolic, and the Imaginary are introduced as three coupled systems, and value is specified as a phenomenon emergent on the dynamics of the coupled whole rather than as a state of any single register. This construal, developed below, determines the sense in which value is at once symbolisable and never exhausted by symbolisation, and thereby corrects a misconstrual to which the description is otherwise liable.

\subsection{The three registers}
\label{sec:rsi}

Lacan distinguishes three registers, or orders, in which psychic and relational reality is structured \citep{evans1996}. The \emph{Imaginary} is the register of the image, of identification and the specular relation, instituted in the mirror stage, in which the subject apprehends itself and the other under the aspect of a unified form and a captating likeness \citep{lacan2006}. The \emph{Symbolic} is the register of the signifier, of language, law, and the differential structure within which terms acquire determinacy through their relations to other terms; it is the order of speech and the social, and the order in which exchange, nomination, and the cashing of a term against others are possible \citep{lacan2006}. The \emph{Real} is the register of what resists symbolisation absolutely, what falls outside the signifier and returns always to the same place; it is approached through the concepts of \emph{das Ding} and \emph{objet a}, and is bound to \emph{jouissance}, the enjoyment beyond the pleasure principle.

Two concepts of the Real require distinction, since the analysis of the following sections turns on it. \emph{Das Ding}, the Thing, is the term of the early account in the \emph{Ethics of Psychoanalysis}: the absolute Other of the subject, an originary void at the centre of the Real around which the subject's economy is organised and which no representation reaches \citep{lacan1992}. \emph{Objet a}, the object-cause of desire, is the term of the later account: not the object at which desire aims but the object that causes it, a remainder that falls from the field of the Other and functions as the cause of desire rather than its goal \citep{lacan1978sem11,fink1995}. The relation between the two is one of development rather than opposition: \emph{objet a} may be construed as the objectalisation, within the economy of desire, of the void that \emph{das Ding} names in the Real. \emph{Jouissance} is the enjoyment associated with this register, transgressive of the pleasure principle, at once sought and unbearable, and it is the quantity that the later argument identifies as the carrier of the relational phase \citep{lacan1998,fink1995}.

\subsection{The registers as coupled systems}
\label{sec:coupled}

In its late development Lacan presents the three registers as knotted, in the figure of the Borromean knot: three rings so linked that the severing of any one releases the other two \citep{lacan1998,evans1996}. The present paper retains the interdependence the knot expresses but reconstrues it dynamically. The static topology of the knot is replaced by a dynamics of coupling: the Real, the Symbolic, and the Imaginary are treated as three coupled systems, each with its own internal dynamics and its own state, interacting through coupling terms by which the state of each conditions the evolution of the others.

\begin{claim}[The registers as coupled dynamical systems]
The three registers are construed as three coupled dynamical systems. Each register possesses an internal dynamics and a state; the registers interact through coupling terms, such that the state of each conditions the evolution of the others. The interdependence Lacan expresses in the static figure of the Borromean knot, that the severing of any one ring releases the others, is reconstrued as the dynamical condition that the suppression of any one coupling alters the dynamics of the coupled whole. The construal is dynamical rather than topological: the registers are not three rings in a fixed linkage but three systems in continuous mutual conditioning.
\end{claim}

\subsection{Value as an emergent of the coupled dynamics}
\label{sec:emergence}

Value is specified, on this construal, as a phenomenon emergent on the dynamics of the coupled whole. It is not a state of the Imaginary alone, not a state of the Symbolic alone, and not a state of the Real alone; it is a collective mode of the three coupled systems, arising on their interaction and not localisable in any one of them.

\begin{claim}[Value as an emergent collective mode of the coupled registers]
Value is not a state of any single register but a phenomenon emergent on the dynamics of the three coupled registers. Formally, where each register is a dynamical system and the registers are coupled, value corresponds to a collective observable of the coupled system, an order parameter that is non-zero only under non-vanishing coupling among the registers and that vanishes when any one coupling is suppressed. Value has accordingly an imaginary aspect (the idealised image, the captating form, the value invested in the likeness), a symbolic aspect (the nameable, exchangeable, cashable value), and a real aspect (the value organised around the void, bound to jouissance, irreducible to the signifier); but value is identical to none of these aspects and is exhausted by none. It is the collective mode in which the three are coupled.
\end{claim}

\noindent
The consequence that governs the subsequent analysis is the following. Since value is an emergent of the coupled dynamics, its reduction to any single register, the suppression of a coupling, extinguishes it. The reduction of value to the Imaginary is its collapse into idealisation, the vertiginous investment in a form that exhausts itself; the reduction to the Symbolic is its collapse into pure exchange, the cashing of the value as a settled token severed from the Real; the reduction to the Real is its collapse into an unsymbolisable remainder that cannot be held or shared. Each reduction suppresses a coupling and extinguishes the emergent. This supplies a dynamical account of a thesis the prior papers asserted: the settling of value into pure symbol, its severance from the Real, does not merely impoverish the value but extinguishes it, value being an emergent that the severance destroys.

This construal corrects a misconstrual to which the phenomenological description is otherwise liable, and which is addressed at \xref{sec:nameless-corrected}. Value is not the unsymbolisable Real, the term that naming would kill; value is the emergent of the coupling of the Real with the Symbolic and the Imaginary, and is therefore both symbolisable and inexhaustible by symbolisation. Two persons do symbolise value: they name it, they make private signs for it, they bring it into the Symbolic. The question is not whether value is symbolised but how. A symbolisation that sustains the coupling of the Symbolic with the Real, the poetic symbolisation of the later analysis, sustains the emergent; a symbolisation that severs the Symbolic from the Real, the settling symbolisation, suppresses the coupling and extinguishes it. The transversality of value across the three registers, its being an emergent of their coupled dynamics rather than a resident of any one, is the structure the following section's datum exhibits, and the structure to which the account of language in this series is addressed.
